\begin{abstract}
The sidewall angle of etched silicon waveguides changes the effective index of every guided mode, yet the optical
test structures used to monitor fabrication do not report it. We analyze which test structures can determine the
angle, using full-vector finite-element mode solutions and the Cram\'er--Rao lower bound (CRLB) of a joint fit of
ring and interferometer readouts. For an oxide-clad strip the mode is mirror-symmetric about mid-height, so to first
order a wall tilt $\theta$ about the bottom edge is exactly a width change of $t\tan\theta$, \qty{3.84}{nm} per degree
for a \qty{220}{nm} film: sets of strips carry no angle information at any size. A partially etched rib breaks the
symmetry, with an angle/width ratio of \numrange{0.75}{0.87} instead of 1, but a rib made by a separate etch brings
its own depth and angle, and one rib width then leaves the strip angle unidentifiable. Two widths of the same rib
resolve it. With strip rings at five widths and ribs of \qty{300}{nm} and \qty{500}{nm} on a \qty{50}{nm} slab, the
CRLB on the strip sidewall angle is \qty{0.39}{\degree} for a \qty{1}{pm} resonance-wavelength precision, with etch
depth, rib angle and both material indices unknown; the best eight-structure set reaches about
\qtyrange{0.12}{0.14}{\degree}, with a \qty{15}{\percent} uncertainty from derivative accuracy. Nonlinear fits of
2000 simulated noisy measurement sets reproduce these bounds to within \qty{5}{\percent}. Treating a \qty{1}{\degree}
tilt as a width error mispredicts a rib ring resonance by \qty{711}{pm}. The results are from simulations only.
\end{abstract}

\begin{IEEEkeywords}
Silicon photonics, sidewall angle, process monitoring, test structures, Cram\'er--Rao bound, ring resonators.
\end{IEEEkeywords}
